% ============================================================
% v2/sections_v2/discussion.tex  考察（短縮・再構成版，査読の指摘を受けて改訂 2026-10-27 17:00）
% 数字は sections_v2/results.tex と同じ出典。新しい数字は sections/discussion.tex（照合済み）から。
% ============================================================
\section{考察}\label{sec:discussion}

\subsection{何が新しく，何が分かったか}\label{sec:disc-what}
文献\cite{noguchi2022superposition}は，自他で重みを共有する予測学習から他者の位置・視点・運動の表現が生まれ，他者の観察から自分の運動に当たる信号が作られることを示した．
自分側で学んだ感覚運動の予測が他者側に使われうること自体は，そこですでに示されている．
本研究で新しく分かったのは，次の三点である．

第一に，与えた価値の信号に従って動いた経験から，報酬を用いない予測学習で価値に従う行動の見込みが生まれ，その見込みが，段階Iでは零しか入らない他者側に与えた快・不快にも反映された．

第二に，この転用は重みの共有だけでは生じなかった．
位置・運動・視界の系列を同じにして価値の信号と行動の対応だけを壊すと，自分側の快・不快への依存と他者側の符号に従う分は0前後になり，G5b も大きく弱まった．
この対照がなければ，経験80\%条件の効果は，赤の角に多くいた経験や，位置の手がかりとしての快・不快の大きさでも説明できた．
識別できたのは価値の信号と行動の一貫した対応の寄与であり，価値に従って動いた経験のあらゆる効果ではない．
G5b が経験なし条件より大きく残ったことは，赤への癖などの残りを含む．

第三に，学習後に写す比較から，転用は快・不快の入口の重みだけでは説明できず，自分の経験で学習された LSTM のほかの重みや計算にも依存する可能性が示された．
ただし，回復したのは視覚予測器を通らない量だけで，どの重みや計算が運ばれるかは特定していない．

予測に役立つ入力が使われるようになること自体は，予測学習の一般的な性質である．
本研究の知見は，その使い方が自分の行動の経験から学ばれて他者側の予測にも反映されること，そしてそれに要る条件の一つ（価値の信号と行動の一貫した対応）を，位置の分布をそろえた対照で識別したことにある．
これは，自分の意思決定の仕組みを他者の行動の仮説を試すために走らせるというシミュレーション説の考え\cite{gordon2021simulation}について，何がそろえば価値に関わる他者予測への転用が起きるかを，一つの構成の中で具体化したものと位置づけられる．
他の構成でも同じ条件が要るかは示していない．

\subsection{主張の強さと，前提を満たさない条件への依存}\label{sec:disc-strength}
経験80\%条件のモデルは，自己位置の追跡の事前基準を6組とも満たさなかった（共通テストで 0.70〜0.75，自分の学習データと同じ作りのテストでは 0.99）．
主張の強さは，これに応じて次のように分けられる．

前提を満たした経験50\%条件で支持されるのは，価値に従う行動の見込みが経験によって生まれること（6組すべてで経験なし条件との差）と，他者側の符号応答 G5b が事前登録した判定に合格したこと（本番の3組中2組）である．
ただし経験50\%条件の G5b は組をまたぐと差が0を含み（表\ref{tab:summary}），転用は弱く，組によってそろわない．

はっきりした転用，成立条件の切り分け（符号無作為条件，写す比較），他者の規則の確率を変えた比較は，すべて経験80\%条件に依っている．
このうち符号無作為条件は，経験80\%条件と同じ軌跡で学習し，自己位置の追跡も経験80\%条件に近かった．
したがって，両者の違いは前提の不成立では説明できない．
一方，経験80\%条件のモデルは評価に用いたランダムな歩きの場面に十分に一般化していないので，よく一般化したモデルで同じ大きさの効果が得られるかは分からない．
成立条件についての結論は，「赤の角にとどまる経験を多く積んだモデルで，価値の信号と行動の対応が転用に要った」という範囲で読むべきである．

\subsection{転用は他者の理解を保証しない}\label{sec:disc-st}
G5b と符号に従う分が示すのは，自分側で学んだ傾向が他者側の予測に反映されることであり，他者の行動の予測が正確になることではない．
実際，本当の快・不快を入れたときの誤差は経験なし条件と経験50\%条件でほぼ同じで，経験80\%条件ではかえって大きくなる組があった（\ref{sec:res-projection}節）．
自分由来の予測の転用は，他者の適切な理解を必ずしも保証しない．

人の他者理解では，自分を出発点にし，そこから他者に合わせて調整する過程が想定されている\cite{epley2004perspective,tamir2010neural}．
Suzuki ら\cite{suzuki2012learning}の区別でいえば，本研究のモデルは自分の価値づけを流用する経路を持つが，他者の行動の予測の外れから他者の価値を補正する経路を持たない．
他者の規則の確率を自分の学習条件に近づけると本当の快・不快がより役立ったことは，調整を欠いた転用が行動の仕方の違う他者で失敗する例と読むこともできる．
この読みで当てはめられているのは「何を好むか」ではなく「どれだけ強く好みに従うか」である．
しかし，規則の確率を変えると位置の分布も変わり，確率が一致したときに最もよく役立つことも示しておらず，符号無作為条件でも悪化の原因は判定できなかった．
この解釈は可能性にとどまる．
なお，自分の行為の経験が他者の行為の理解を変えること\cite{sommerville2005action}や，幼い子どもが他者に自分の好みを当てはめること\cite{repacholi1997early}と形の似た振る舞いも見られたが，これは計算モデルの振る舞いとの類推であり，人間について何かを示すものではない．

\subsection{なぜ他者の好みは取り出されないか}\label{sec:disc-shortcut}
他者用モジュールは他者側の予測の外れを減らすように学習するので，補正の学習信号そのものは与えられていた．
実際，経験80\%条件の上で学習したモジュールの出力を入れると，他者側の特徴予測の誤差は 0.112〜0.188 で，本当の快・不快を入れた場合（0.376〜0.418）よりずっと小さかった（組1〜4，二つのデータ）．
補正は学ばれたが，他者の価値を直す形ではなく，予測に役立つ位置や方角を伝える形をとった．
これは，学習の目的（予測の外れを減らすこと）はよく満たすが，意図した表現（他者の価値）とは別の手がかりに頼る近道\cite{geirhos2020shortcut}に当たると考えられる（推測であり，直接操作して確かめてはいない）．
本研究の課題では A-2 は一歩ごとに約半分しか規則に従わないので，本当の快・不快は A-2 の動きを部分的にしか決めず，位置や方角を伝える方が予測をずっとよくする．
なお，出力に好みが表れないことは，内部状態に好みの情報がないことを意味しない．

\subsection{限界}\label{sec:limitations}
第一に，環境は小さい．四隅にランドマークを置いた一つのアリーナで，カメラは回転しないので，「赤の方へ動く」ことは「赤の象限の方へ動く」ことと区別しにくい．
第二に，自他の対応は基準エンコーダで与え，価値は人が方角ごとに計算して与えた信号である．
本研究が示したのは，与えた価値の信号とそれに従う行動の規則性を予測学習が利用することまでである．
第三に，前提を満たさない条件への依存は\ref{sec:disc-strength}節で述べたとおりである．
第四に，非共有条件で他者側の働きが消えることは構成からほぼ予想され，写す比較は事前登録外で，予測画像を通る量では確かめられなかった．
第五に，符号無作為条件は行動の履歴そのものを変えていない．快・不快を常に零にした条件の結果は【結果待ち】である．
第六に，独立に学習した組は本番の3組，追加を含めて6組（符号無作為条件と写す比較は3組）と少なく，組ごとのばらつきも大きい．最終5点は互いに相関するチェックポイントで独立した繰り返しではなく，多重比較の補正も行っていない．符号に従う分（組1〜4）は事後に定義した指標である．
第七に，照合の候補は研究者が用意した二択であり，ネットワークが自ら他者の好みを推し量ったとは言えない．

\subsection{今後の課題}\label{sec:future}
\ref{sec:disc-shortcut}節の解釈からは，他者の行動の仕方が自分と同じなら，他者用モジュールは価値を書き出すようになるはずだという予測が出る．
本当の快・不快が役立つようになるという前半は\ref{sec:res-projection}節の比較で支持されたので，後半を確かめるには，A-2 が 0.8 の確率で規則に従うデータで他者用モジュールを学習する必要がある．
また，好きな物に着くとそれが価値を失い，次に好きな物へ向かう環境のように，好みの順位でしか次の行き先が決まらない課題では，位置や方角の近道が使えなくなるので，照合や他者用モジュールが方角を超えて好みを取り出すかを調べられる．
ただし，観察から順位を推す材料が本当にあるかを，学習の前にデータだけで確かめる必要がある．
そのほか，符号無作為条件の組を増やし，自己位置の追跡の前提を満たす条件でも同じ対照を行うこと，写した状態に合わせて視覚予測器を学習し直して予測画像を通る量でも確かめること，他者の行動の予測の外れから他者の価値を補正する経路を加えること，自他の対応と価値の表現を同時に学ぶエンコーダで同じ結果が得られるかを確かめることが挙げられる．
